\section{激活函数}

\subsection{ReLU 及其变体}

\textbf{ReLU}（Rectified Linear Unit）是最广泛使用的激活函数：

$$\text{ReLU}(x) = \max(0, x)$$

\begin{figure}[H]
\centering
\includegraphics[width=0.95\textwidth]{slides-images/slide-028.jpg}
\caption{ReLU 激活函数：$x > 0$ 时输出 $x$，$x \leq 0$ 时输出 0\protect\footnotemark}
\end{figure}
\footnotetext{Slides 第29页。}

ReLU 的\textbf{问题}：当输入为负时，输出恒为 0，梯度也为 0，导致所谓的``死神经元''——一旦某个神经元的输入持续为负，它将永远无法恢复。

为此，出现了多种改进变体：

\begin{figure}[H]
\centering
\includegraphics[width=0.95\textwidth]{slides-images/slide-029.jpg}
\caption{ReLU 的变体：Leaky ReLU、ELU、GELU、SiLU/Swish\protect\footnotemark}
\end{figure}
\footnotetext{Slides 第30页。}

\begin{itemize}
    \item \textbf{Leaky ReLU}：$f(x) = \max(0.01x, x)$，负半轴给一个小斜率，避免死神经元
    \item \textbf{ELU}（Exponential Linear Unit）：负半轴使用指数函数，具有更好的零中心性
    \item \textbf{GELU}（Gaussian Error Linear Unit）：在 Transformer 架构中广泛使用
    \item \textbf{SiLU/Swish}：$f(x) = x \cdot \sigma(x)$，Google 在 EfficientNet 等模型中采用
\end{itemize}

\subsection{Sigmoid 和 Tanh}

\begin{figure}[H]
\centering
\includegraphics[width=0.95\textwidth]{slides-images/slide-030.jpg}
\caption{Sigmoid 和 Tanh 激活函数：将输出压缩到有限范围内\protect\footnotemark}
\end{figure}
\footnotetext{Slides 第31页。}

\textbf{Sigmoid}：$\sigma(x) = \frac{1}{1 + e^{-x}}$，输出范围 $(0, 1)$

\textbf{Tanh}：$\tanh(x) = \frac{e^x - e^{-x}}{e^x + e^{-x}}$，输出范围 $(-1, 1)$

\begin{warningbox}{Sigmoid 和 Tanh 的梯度消失问题}
这两个函数将值压缩到很窄的范围内。当输入值的绝对值很大时，函数曲线几乎是平的，导数接近零。在深层网络中，这些接近零的梯度经过链式法则逐层相乘，最终变得极其微小——这就是\textbf{梯度消失}（Vanishing Gradient）问题。因此，Sigmoid 和 Tanh 通常\textbf{不用在网络的中间层}，而只用在输出层（如二分类的 Sigmoid 输出）。
\end{warningbox}

\subsection{激活函数的选择}

\begin{knowledgebox}{激活函数选择指南}
\begin{itemize}
    \item \textbf{默认选择}：ReLU，简单高效，在大多数场景下表现良好
    \item \textbf{CNN 架构}：ReLU 或 SiLU/Swish
    \item \textbf{Transformer 架构}：GELU 是标准选择
    \item \textbf{输出层}：Sigmoid（二分类）、Softmax（多分类）、无激活（回归）
    \item 所有激活函数的共同要求：\textbf{非线性}且\textbf{可微}
\end{itemize}
选择激活函数更多是经验性的——通常沿用同类型架构已验证有效的选择。
\end{knowledgebox}

\subsection{本章小结}

激活函数是神经网络的核心组件，为网络提供非线性能力。ReLU 是最常用的默认选择，但存在死神经元问题。GELU 和 SiLU 是更现代的替代方案。Sigmoid 和 Tanh 因梯度消失问题通常不用于中间层。

%% ========== 网络实现 ==========

